% figA10_finetune.tex — Figure A10: Task fine-tuning blind spot hierarchy
% % figA10_finetune.tex — Figure A10: Task fine-tuning blind spot hierarchy
% % figA10_finetune.tex — Figure A10: Task fine-tuning blind spot hierarchy
% % figA10_finetune.tex — Figure A10: Task fine-tuning blind spot hierarchy
% \input{figures/figA10_finetune}

\begin{figure*}[t]
\centering
\begin{subfigure}[b]{0.50\textwidth}
\begin{tikzpicture}
\begin{axis}[
  pmh bar,
  symbolic x coords={Pretrained,{ERM fine-tuned},{PMH fine-tuned}},
  xtick=data,
  xticklabel style={font=\footnotesize, align=center},
  ylabel={Paraphrase drift (CLS displacement)},
  ymin=0, ymax=0.056,
  width=0.92\linewidth, height=5.5cm,
  bar width=9pt,
  nodes near coords,
  nodes near coords style={font=\tiny, anchor=south},
  point meta=rawy,
  every node near coord/.append style={
    /pgf/number format/fixed, /pgf/number format/precision=5},
]
\addplot[fill=colB0, draw=none]
  coordinates {(Pretrained,0.02435)};
\addplot[fill=colVAT, draw=none]
  coordinates {({ERM fine-tuned},0.03751)};
\addplot[fill=colPMH, draw=none]
  coordinates {({PMH fine-tuned},0.00333)};
\node[font=\tiny\bfseries, color=colRed, align=center]
  at (axis cs:{ERM fine-tuned},0.0510) {ERM fine-tune\\blind spot $\uparrow$};
\node[font=\tiny\bfseries, color=colGreen, align=center]
  at (axis cs:{PMH fine-tuned},0.0510) {PMH repairs\\$11\times$ $\downarrow$};
\end{axis}
\end{tikzpicture}
\caption{Predicted ordering ERM $>$ pretrained $>$ PMH confirmed.}
\end{subfigure}
\hfill
\begin{subfigure}[b]{0.46\textwidth}
\centering
\begin{tikzpicture}[
  every node/.style={font=\small},
  box/.style={draw, rounded corners=3pt, inner sep=6pt, line width=0.6pt,
              minimum width=3.6cm, align=center},
  arr/.style={-{Latex[length=3.5pt]}, line width=0.8pt},
]
\node[box, draw=colB0!50, fill=colB0!8] (pre) at (0,3.4)
  {\footnotesize\textbf{Pretrained}\\[2pt]
   \tiny drift $= 0.0244$\\
   \tiny blind spot ratio $= 0.765$};
\node[box, draw=colRed!50, fill=colRed!8] (erm) at (0,1.7)
  {\footnotesize\textbf{ERM fine-tuned}\\[2pt]
   \tiny drift $= 0.0375$\\
   \tiny blind spot ratio $= 0.681$};
\node[box, draw=colGreen!60, fill=colGreen!8] (pmh) at (0,0.0)
  {\footnotesize\textbf{PMH fine-tuned}\\[2pt]
   \tiny drift $= 0.0033$\\
   \tiny blind spot ratio $= 0.633$};

\draw[arr, colRed] (pre) -- (erm)
  node[midway, anchor=west, xshift=2.0cm,
       font=\tiny, color=colRed, align=left]
    {task fine-tuning\\blind spot $\uparrow$};
\draw[arr, colGreen] (erm) -- (pmh)
  node[midway, anchor=west, xshift=2.0cm,
       font=\tiny, color=colGreen, align=left]
    {$+$ PMH term\\drift $11\times$ $\downarrow$};

\node[font=\tiny, gray, align=center] at (0,-0.65)
  {Predicted: ERM $>$ pretrained $>$ PMH\ $\checkmark$};
\end{tikzpicture}
\caption{Blind spot hierarchy.}
\end{subfigure}
\caption{\textbf{Task fine-tuning worsens the geometric blind spot; PMH repairs it.}
\emph{(a)} Paraphrase drift follows the predicted theoretical hierarchy:
ERM fine-tuning (0.0375) $>$ pretrained baseline (0.0244) $>$ PMH fine-tuning (0.0033).
Task-specific ERM amplifies nuisance encoding (increases effective $\rho$);
PMH suppresses it $11\times$.
\emph{(b)} The blind spot ratio decreases monotonically through the training hierarchy: pretrained (0.765) $\to$ ERM fine-tuned (0.681) $\to$ PMH (0.633).
This is not merely compensating for fine-tuning noise: PMH repairs the geometric structure
predicted by Theorem~1, confirming the fix generalises to the full supervised fine-tuning
hierarchy present in modern pre-trained language model pipelines.}
\label{fig:A10}
\end{figure*}


\begin{figure*}[t]
\centering
\begin{subfigure}[b]{0.50\textwidth}
\begin{tikzpicture}
\begin{axis}[
  pmh bar,
  symbolic x coords={Pretrained,{ERM fine-tuned},{PMH fine-tuned}},
  xtick=data,
  xticklabel style={font=\footnotesize, align=center},
  ylabel={Paraphrase drift (CLS displacement)},
  ymin=0, ymax=0.056,
  width=0.92\linewidth, height=5.5cm,
  bar width=9pt,
  nodes near coords,
  nodes near coords style={font=\tiny, anchor=south},
  point meta=rawy,
  every node near coord/.append style={
    /pgf/number format/fixed, /pgf/number format/precision=5},
]
\addplot[fill=colB0, draw=none]
  coordinates {(Pretrained,0.02435)};
\addplot[fill=colVAT, draw=none]
  coordinates {({ERM fine-tuned},0.03751)};
\addplot[fill=colPMH, draw=none]
  coordinates {({PMH fine-tuned},0.00333)};
\node[font=\tiny\bfseries, color=colRed, align=center]
  at (axis cs:{ERM fine-tuned},0.0510) {ERM fine-tune\\blind spot $\uparrow$};
\node[font=\tiny\bfseries, color=colGreen, align=center]
  at (axis cs:{PMH fine-tuned},0.0510) {PMH repairs\\$11\times$ $\downarrow$};
\end{axis}
\end{tikzpicture}
\caption{Predicted ordering ERM $>$ pretrained $>$ PMH confirmed.}
\end{subfigure}
\hfill
\begin{subfigure}[b]{0.46\textwidth}
\centering
\begin{tikzpicture}[
  every node/.style={font=\small},
  box/.style={draw, rounded corners=3pt, inner sep=6pt, line width=0.6pt,
              minimum width=3.6cm, align=center},
  arr/.style={-{Latex[length=3.5pt]}, line width=0.8pt},
]
\node[box, draw=colB0!50, fill=colB0!8] (pre) at (0,3.4)
  {\footnotesize\textbf{Pretrained}\\[2pt]
   \tiny drift $= 0.0244$\\
   \tiny blind spot ratio $= 0.765$};
\node[box, draw=colRed!50, fill=colRed!8] (erm) at (0,1.7)
  {\footnotesize\textbf{ERM fine-tuned}\\[2pt]
   \tiny drift $= 0.0375$\\
   \tiny blind spot ratio $= 0.681$};
\node[box, draw=colGreen!60, fill=colGreen!8] (pmh) at (0,0.0)
  {\footnotesize\textbf{PMH fine-tuned}\\[2pt]
   \tiny drift $= 0.0033$\\
   \tiny blind spot ratio $= 0.633$};

\draw[arr, colRed] (pre) -- (erm)
  node[midway, anchor=west, xshift=2.0cm,
       font=\tiny, color=colRed, align=left]
    {task fine-tuning\\blind spot $\uparrow$};
\draw[arr, colGreen] (erm) -- (pmh)
  node[midway, anchor=west, xshift=2.0cm,
       font=\tiny, color=colGreen, align=left]
    {$+$ PMH term\\drift $11\times$ $\downarrow$};

\node[font=\tiny, gray, align=center] at (0,-0.65)
  {Predicted: ERM $>$ pretrained $>$ PMH\ $\checkmark$};
\end{tikzpicture}
\caption{Blind spot hierarchy.}
\end{subfigure}
\caption{\textbf{Task fine-tuning worsens the geometric blind spot; PMH repairs it.}
\emph{(a)} Paraphrase drift follows the predicted theoretical hierarchy:
ERM fine-tuning (0.0375) $>$ pretrained baseline (0.0244) $>$ PMH fine-tuning (0.0033).
Task-specific ERM amplifies nuisance encoding (increases effective $\rho$);
PMH suppresses it $11\times$.
\emph{(b)} The blind spot ratio decreases monotonically through the training hierarchy: pretrained (0.765) $\to$ ERM fine-tuned (0.681) $\to$ PMH (0.633).
This is not merely compensating for fine-tuning noise: PMH repairs the geometric structure
predicted by Theorem~1, confirming the fix generalises to the full supervised fine-tuning
hierarchy present in modern pre-trained language model pipelines.}
\label{fig:A10}
\end{figure*}


\begin{figure*}[t]
\centering
\begin{subfigure}[b]{0.50\textwidth}
\begin{tikzpicture}
\begin{axis}[
  pmh bar,
  symbolic x coords={Pretrained,{ERM fine-tuned},{PMH fine-tuned}},
  xtick=data,
  xticklabel style={font=\footnotesize, align=center},
  ylabel={Paraphrase drift (CLS displacement)},
  ymin=0, ymax=0.056,
  width=0.92\linewidth, height=5.5cm,
  bar width=9pt,
  nodes near coords,
  nodes near coords style={font=\tiny, anchor=south},
  point meta=rawy,
  every node near coord/.append style={
    /pgf/number format/fixed, /pgf/number format/precision=5},
]
\addplot[fill=colB0, draw=none]
  coordinates {(Pretrained,0.02435)};
\addplot[fill=colVAT, draw=none]
  coordinates {({ERM fine-tuned},0.03751)};
\addplot[fill=colPMH, draw=none]
  coordinates {({PMH fine-tuned},0.00333)};
\node[font=\tiny\bfseries, color=colRed, align=center]
  at (axis cs:{ERM fine-tuned},0.0510) {ERM fine-tune\\blind spot $\uparrow$};
\node[font=\tiny\bfseries, color=colGreen, align=center]
  at (axis cs:{PMH fine-tuned},0.0510) {PMH repairs\\$11\times$ $\downarrow$};
\end{axis}
\end{tikzpicture}
\caption{Predicted ordering ERM $>$ pretrained $>$ PMH confirmed.}
\end{subfigure}
\hfill
\begin{subfigure}[b]{0.46\textwidth}
\centering
\begin{tikzpicture}[
  every node/.style={font=\small},
  box/.style={draw, rounded corners=3pt, inner sep=6pt, line width=0.6pt,
              minimum width=3.6cm, align=center},
  arr/.style={-{Latex[length=3.5pt]}, line width=0.8pt},
]
\node[box, draw=colB0!50, fill=colB0!8] (pre) at (0,3.4)
  {\footnotesize\textbf{Pretrained}\\[2pt]
   \tiny drift $= 0.0244$\\
   \tiny blind spot ratio $= 0.765$};
\node[box, draw=colRed!50, fill=colRed!8] (erm) at (0,1.7)
  {\footnotesize\textbf{ERM fine-tuned}\\[2pt]
   \tiny drift $= 0.0375$\\
   \tiny blind spot ratio $= 0.681$};
\node[box, draw=colGreen!60, fill=colGreen!8] (pmh) at (0,0.0)
  {\footnotesize\textbf{PMH fine-tuned}\\[2pt]
   \tiny drift $= 0.0033$\\
   \tiny blind spot ratio $= 0.633$};

\draw[arr, colRed] (pre) -- (erm)
  node[midway, anchor=west, xshift=2.0cm,
       font=\tiny, color=colRed, align=left]
    {task fine-tuning\\blind spot $\uparrow$};
\draw[arr, colGreen] (erm) -- (pmh)
  node[midway, anchor=west, xshift=2.0cm,
       font=\tiny, color=colGreen, align=left]
    {$+$ PMH term\\drift $11\times$ $\downarrow$};

\node[font=\tiny, gray, align=center] at (0,-0.65)
  {Predicted: ERM $>$ pretrained $>$ PMH\ $\checkmark$};
\end{tikzpicture}
\caption{Blind spot hierarchy.}
\end{subfigure}
\caption{\textbf{Task fine-tuning worsens the geometric blind spot; PMH repairs it.}
\emph{(a)} Paraphrase drift follows the predicted theoretical hierarchy:
ERM fine-tuning (0.0375) $>$ pretrained baseline (0.0244) $>$ PMH fine-tuning (0.0033).
Task-specific ERM amplifies nuisance encoding (increases effective $\rho$);
PMH suppresses it $11\times$.
\emph{(b)} The blind spot ratio decreases monotonically through the training hierarchy: pretrained (0.765) $\to$ ERM fine-tuned (0.681) $\to$ PMH (0.633).
This is not merely compensating for fine-tuning noise: PMH repairs the geometric structure
predicted by Theorem~1, confirming the fix generalises to the full supervised fine-tuning
hierarchy present in modern pre-trained language model pipelines.}
\label{fig:A10}
\end{figure*}


\begin{figure*}[t]
\centering
\begin{subfigure}[b]{0.50\textwidth}
\begin{tikzpicture}
\begin{axis}[
  pmh bar,
  symbolic x coords={Pretrained,{ERM fine-tuned},{PMH fine-tuned}},
  xtick=data,
  xticklabel style={font=\footnotesize, align=center},
  ylabel={Paraphrase drift (CLS displacement)},
  ymin=0, ymax=0.056,
  width=0.92\linewidth, height=5.5cm,
  bar width=9pt,
  nodes near coords,
  nodes near coords style={font=\tiny, anchor=south},
  point meta=rawy,
  every node near coord/.append style={
    /pgf/number format/fixed, /pgf/number format/precision=5},
]
\addplot[fill=colB0, draw=none]
  coordinates {(Pretrained,0.02435)};
\addplot[fill=colVAT, draw=none]
  coordinates {({ERM fine-tuned},0.03751)};
\addplot[fill=colPMH, draw=none]
  coordinates {({PMH fine-tuned},0.00333)};
\node[font=\tiny\bfseries, color=colRed, align=center]
  at (axis cs:{ERM fine-tuned},0.0510) {ERM fine-tune\\blind spot $\uparrow$};
\node[font=\tiny\bfseries, color=colGreen, align=center]
  at (axis cs:{PMH fine-tuned},0.0510) {PMH repairs\\$11\times$ $\downarrow$};
\end{axis}
\end{tikzpicture}
\caption{Predicted ordering ERM $>$ pretrained $>$ PMH confirmed.}
\end{subfigure}
\hfill
\begin{subfigure}[b]{0.46\textwidth}
\centering
\begin{tikzpicture}[
  every node/.style={font=\small},
  box/.style={draw, rounded corners=3pt, inner sep=6pt, line width=0.6pt,
              minimum width=3.6cm, align=center},
  arr/.style={-{Latex[length=3.5pt]}, line width=0.8pt},
]
\node[box, draw=colB0!50, fill=colB0!8] (pre) at (0,3.4)
  {\footnotesize\textbf{Pretrained}\\[2pt]
   \tiny drift $= 0.0244$\\
   \tiny blind spot ratio $= 0.765$};
\node[box, draw=colRed!50, fill=colRed!8] (erm) at (0,1.7)
  {\footnotesize\textbf{ERM fine-tuned}\\[2pt]
   \tiny drift $= 0.0375$\\
   \tiny blind spot ratio $= 0.681$};
\node[box, draw=colGreen!60, fill=colGreen!8] (pmh) at (0,0.0)
  {\footnotesize\textbf{PMH fine-tuned}\\[2pt]
   \tiny drift $= 0.0033$\\
   \tiny blind spot ratio $= 0.633$};

\draw[arr, colRed] (pre) -- (erm)
  node[midway, anchor=west, xshift=2.0cm,
       font=\tiny, color=colRed, align=left]
    {task fine-tuning\\blind spot $\uparrow$};
\draw[arr, colGreen] (erm) -- (pmh)
  node[midway, anchor=west, xshift=2.0cm,
       font=\tiny, color=colGreen, align=left]
    {$+$ PMH term\\drift $11\times$ $\downarrow$};

\node[font=\tiny, gray, align=center] at (0,-0.65)
  {Predicted: ERM $>$ pretrained $>$ PMH\ $\checkmark$};
\end{tikzpicture}
\caption{Blind spot hierarchy.}
\end{subfigure}
\caption{\textbf{Task fine-tuning worsens the geometric blind spot; PMH repairs it.}
\emph{(a)} Paraphrase drift follows the predicted theoretical hierarchy:
ERM fine-tuning (0.0375) $>$ pretrained baseline (0.0244) $>$ PMH fine-tuning (0.0033).
Task-specific ERM amplifies nuisance encoding (increases effective $\rho$);
PMH suppresses it $11\times$.
\emph{(b)} The blind spot ratio decreases monotonically through the training hierarchy: pretrained (0.765) $\to$ ERM fine-tuned (0.681) $\to$ PMH (0.633).
This is not merely compensating for fine-tuning noise: PMH repairs the geometric structure
predicted by Theorem~1, confirming the fix generalises to the full supervised fine-tuning
hierarchy present in modern pre-trained language model pipelines.}
\label{fig:A10}
\end{figure*}
